\PassOptionsToPackage{usenames,x11names, dvipsnames, svgnames}{xcolor}
\documentclass[9pt,twocolumn]{IEEEtran}
% Use the lineno option to display guide line numbers if required.
\usepackage{orcidlink}   % <— add this
% \input{preamble_ieee.tex}
\usepackage{tikz}
\usepackage{pgfplots}
\pgfplotsset{compat=1.18}
%\usepackage{pdfpages}
%\usepackage{shellesc}
%\ShellEscape{pdflatex si}
\usetikzlibrary{shapes,calc,shadows,fadings,arrows,decorations.pathreplacing,automata,positioning}
%\usepackage{hyperref}
\usetikzlibrary{external}
\tikzexternalize[prefix=./Figures/External/]% activate
\tikzexternaldisable 
\usetikzlibrary{decorations.text}
\usepgfplotslibrary{colorbrewer} 
\usepgfplotslibrary{statistics}
\usepgfplotslibrary{fillbetween}
%\usepackage[linesnumbered,ruled,vlined]{algorithm2e}
\usepackage[algo2e,ruled,vlined]{algorithm2e}

\usepackage{mathtools}
\usepackage{amssymb,amsfonts,amsmath,amsthm}

\usepackage{flushend}

\newtheorem{theorem}{Theorem}
\newtheorem{lemma}{Lemma}
\newtheorem{corollary}{Corollary}
\newtheorem{definition}{Definition}

\newtheorem{exmpl}{Example}
\newtheorem{rem}{Remark}
\newtheorem{notn}{Notation}
\usepackage{xspace}
\newcommand{\abs}[1]{\left\lvert #1 \right\rvert}%

\newcommand{\set}[1]{\left\{ #1 \right\}}
\newcommand{\paren}[1]{\left( #1 \right)}
\newcommand{\bracket}[1]{\left[ #1 \right]}
% \newcommand{\norm}[1]{\left\Vert #1 \right\Vert}
\newcommand{\nrm}[1]{\left\llbracket{#1}\right\rrbracket}
\newcommand{\parenBar}[2]{\paren{#1\,{\left\Vert\,#2\right.}}}
\newcommand{\parenBarl}[2]{\paren{\left.#1\,\right\Vert\,#2}}
\newcommand{\ie}{$i.e.$\xspace}
\newcommand{\addcitation}{\textcolor{black!50!red}{\textbf{ADD CITATION}}}
\newcommand{\subtochange}[1]{{\color{black!50!green}{#1}}}
\newcommand{\tobecompleted}{{\color{black!50!red}TO BE COMPLETED.}}

\DeclareMathOperator*{\argmax}{argmax}
\DeclareMathOperator*{\argmin}{arg\,min}
\DeclareMathOperator*{\expect}{\mathbf{E}}
\DeclareMathOperator*{\var}{\mathbf{Var}}

\newcommand{\pIn}{\mathscr{P}_{\textrm{in}}}
\newcommand{\pOut}{\mathscr{P}_{\textrm{out}}}
\newcommand{\aIn}[1][\Sigma]{#1_{\textrm{in}}}
\newcommand{\aOut}[1][\Sigma]{#1_{\textrm{out}}}
\newcommand{\xin}[1]{#1_{\textrm{in}}}
\newcommand{\xout}[1]{#1_{\textrm{out}}}

\newcommand{\R}{\mathbb{R}} % Set of real numbers
\newcommand{\F}[1][]{\mathcal{F}_{#1}}
\newcommand{\SR}{\mathcal{S}} % Semiring of sets
\newcommand{\RR}{\mathcal{R}} % Ring of sets
\newcommand{\N}{\mathbb{N}} % Set of natural numbers (0 included)

\newcommand{\pitilde}{\widetilde{\pi}}
\newcommand{\Pitilde}{\widetilde{\Pi}}

\newcommand{\Pp}[1][n]{\mathscr{P}^+_{#1}}
%% ## Equation Space Control---------------------------
\def\EQSP{2pt}
\newcommand{\mltlne}[2][\EQSP]{\begingroup\setlength\abovedisplayskip{#1}\setlength\belowdisplayskip{#1}\begin{equation}\begin{multlined} #2 \end{multlined}\end{equation}\endgroup\noindent}
\newcommand{\cgather}[2][\EQSP]{\begingroup\setlength\abovedisplayskip{#1}\setlength\belowdisplayskip{#1}\begin{gather} #2 \end{gather}\endgroup\noindent}
\newcommand{\cgathers}[2][\EQSP]{\begingroup\setlength\abovedisplayskip{#1}\setlength\belowdisplayskip{#1}\begin{gather*} #2 \end{gather*}\endgroup\noindent}
\newcommand{\calign}[2][\EQSP]{\begingroup\setlength\abovedisplayskip{#1}\setlength\belowdisplayskip{#1}\begin{align} #2 \end{align}\endgroup\noindent}
\newcommand{\caligns}[2][\EQSP]{\begingroup\setlength\abovedisplayskip{#1}\setlength\belowdisplayskip{#1}\begin{align*} #2 \end{align*}\endgroup\noindent}
\newcommand{\mnp}[2]{\begin{minipage}{#1}#2\end{minipage}} 
\newif\iftikzX
\tikzXtrue
\tikzXfalse 
\newif\ifFIGS  
\FIGSfalse  
\FIGStrue

\def\EXTENDEDDATA{Extended Data\xspace}
\def\SUPPLEMENTARY{Supplementary\xspace}
\def\Methods{Methods\xspace}
\newcounter{Dcounter}
\setcounter{Dcounter}{1}
\newcommand{\DQS}[1]{
    \marginpar{
      \tikzexternaldisable 
      \tikz{
        \node[
          rounded corners=5pt,
          draw=none,
          thick,
          fill=black!10,
          font=\sffamily\fontsize{7}{8}\selectfont
        ]{\mnp{.3in} {\color{Red3}\raggedright  \#\theDcounter.~#1}}; 
      }
    }
  \stepcounter{Dcounter}\xspace
}
\renewcommand{\baselinestretch}{.935}

\newcommand{\K}{K}
\newcommand{\Hamming}{\mathrm{H}}
\newcommand{\SIapp}{SI Appendix}

\usepackage{iftex}

\ifPDFTeX
  \usepackage[T1]{fontenc}
  \usepackage{lmodern}
\else
  \usepackage{fontspec}
  \setmainfont{Latin Modern Roman}
\fi

\def\TITLE{Escape from Typicality: Why Evolution Operates Near One Mutation per Genome per Generation}

\title{Author Response: \TITLE \vspace{-20pt}
}
\newcommand{\issue}[1]{\par\noindent{\bf\sffamily\color{IndianRed2}\colorbox{IndianRed1}{\small \color{white}$\textbullet$\textbf{Issue:}} #1}\par\par\vspace{.3em}}
\newcommand{\authorresponse}[1]{\noindent{\color{MediumBlue}\textbf{Author response:} #1}\par\vspace{1em}}

\begin{document}  

\maketitle


\issue{Editor: “the current manuscript lacks substantial discussion in the context of evolutionary biology and is severely under-referenced.”}
\authorresponse{Addressed. Added explicit evolutionary-biology framing in the Introduction and Discussion, and added biology-facing references on evolvability and mutation-rate evolution. Now we have 13 references instead of the original 8.}

\issue{Editor: “In particular, I am quite surprised that you do not at all discuss evolution of evolvability.”}
\authorresponse{Addressed. Added explicit discussion of evolvability and cited evolvability references.}

\issue{Editor: “I would like to invite you to submit a revision in which the discussion of your results is much more substantially and explicitly connected to concepts of evolutionary biology.”}
\authorresponse{Addressed. Added explicit links to mutation-rate evolution, modifiers, evolvability, and drift-barrier arguments in the main text.}

%\issue{“Terminology may obscure the core argument, where ‘atypical’, ‘exceptional’, and ‘rare’ mutations seem synonymous with ‘randomness deficiency’...”}
%\authorresponse{Partly addressed. Main text now emphasizes mutation-atypicality relative to the mutation kernel. To Do: further terminology simplification in final trimming.}

\issue{Reviewer 1: “It is unclear how to interpret the H3N2 Influenza A data in Figure 1(b). Is it consistently incompressible or showing an interesting dynamic change in 2021-2?”}
\authorresponse{Addressed. Clarification added in Supplementray Methods to point out that the variation only looks large due to the scale of the plot, and the variation is with 1.92 to 1.94 bits per base, which is close to the maximum value of 2.0. }

\issue{Reviewer 1: “Are other flu serotypes excluded shown simply because they were less available, or for some other reason?”}
\authorresponse{Addressed. Comment added to Supplementary Methods that additional species show the same key point that geomes are not very compressible.}

\issue{Reviewer 1: “Is the dip in those two years not informative compared to other data points, which seem more consistent?”}
\authorresponse{Addressed. While there might be some interesting information that these dips reveal, it is  not particularly relevant  to the the points made in this paper.}

\issue{Reviewer 1: “Overall, while the theoretical basis could be simplified to promote accessibility for wider a readership...”}
\authorresponse{Addressed. More discussion added connecting to evolvability and biological framing.}

\issue{Reviewer 2: “The Drake rule is an old observation that mutation rate in prokaryotes is inverse proportional to the genome size and is roughly constant per genome at \textasciitilde{}0.003. This rule has not been found to be exact throughout prokaryotes in the more recent literature (e.g. see Sung et al., PNAS 2012).”}
\authorresponse{Addressed. The revision cites Sung et al. and frames the result as an order-one  regime rather than an exact universal realized value.}

\issue{Reviewer 2: “However, it is absolutely unclear how and why mutation rate would evolve to this genomic value.”%% }
%% \authorresponse{Addressed. Added discussion of mutation-rate modifiers and indirect selection, plus Fig.~1e as an explicit evolutionary illustration.}

%% \issue{Reviewer 2: “
The manuscript lacks any argument on the dynamics of mutation rate modifiers (secondary selection) that may lead to this mutation rate tuning.”}
\authorresponse{Addressed. Added discussion of mutation-rate modifiers and indirect selection, plus Fig.~1e as an explicit evolutionary illustration.}

\issue{Reviewer 2: “It also lacks any data support.”}
\authorresponse{Addressed. Added Fig.~1e as a biological simulation demonstration, while retaining the compression and bounded-complexity panels. It seems that for a conceptual paper such as this, with  length restrictions of a Brief report, the support is adeqaute.}

\issue{Reviewer 2: “In sum, although an interesting read, I do not see that this work brings us closer to understanding of the fundamental principles of mutational variation.”}
\authorresponse{Addressed. The revision now positions the result explicitly as a complementary supply-side principle connected to evolvability and mutation-rate evolution.}






\end{document}
